\chapter{Introduction}
\label{ch:intro}

The rise of large language models (LLMs) has brought with it the voice agent: a system that can
hold a conversation and answer a request, including in demanding fields such as
healthcare. This is the setting for the present project, which was carried out at Tennor,
a company that builds voice agents for healthcare, able to identify patients, answer their
requests, and interact with the software used across the medical sector.

These agents share a design flaw. Because they are built to respond as quickly as
possible, they tend to cut the caller off when a sentence is delivered in several pieces,
for example a spelled-out detail, a telephone number given in blocks, or a pause for
thought. 

A person handles these pauses easily, from the meaning of the words and the
speaker's intonation. A machine that listens only for silence does not. The central
problem of this dissertation is therefore the following: how can the system tell that a
caller has not finished speaking, so that it does not interrupt them, lose information,
and degrade their experience? The decision has to be made in real time after every
silence, and it rests on an uncomfortable trade-off: ending the turn too early talks over
the caller, while waiting too long makes the exchange feel slow.

The title of this dissertation says end-of-utterance; the literature and the benchmark used
here say end-of-turn. They name the same event, the moment a speaker has finished and the
floor is free, and the two are used interchangeably below.

The existing system already separates speech from silence with a voice activity detector
(VAD) and ends the turn after a fixed timeout. It listens to the sound, but not to the
meaning. My idea is to add a semantic signal: a language model reads
the partial transcript and estimates whether the words form a complete utterance. This
signal is fused with the acoustic one into a single probability, that probability is
calibrated, and a stop-or-wait decision is then derived, in the offline evaluation, from an explicit
cost: the price of an interruption set against that of a delay, rather than a simple timer.

Fusing an acoustic signal with a lexical one is not, in itself, new. The contribution of
this project is narrower: a composition of off-the-shelf models, with no joint training and
no large language model in the decision loop; a probability that is validated before it is
trusted; and a decision derived from an explicit cost. The whole
study is grounded in a real case of French medical telephony, with an in-domain dataset
gathered by instrumenting the production system.

How that claim came to be scoped so narrowly is part of the work. When I chose this project
and had it approved, in April 2026, systems fusing an acoustic and a lexical signal existed,
as Chapter~\ref{ch:background} sets out, but none I could find treated its output as a
probability to be calibrated or derived its endpoint from a stated cost. 

The critical review I submitted on 12 June, of a predictive-endpointing paper \citep{zink2024predictive}, ended
by proposing the three things that became this dissertation: the fusion, the calibrated
cost-sensitive decision, and the design that never waits on the transcript. Five days later
LiveKit released a turn detector and an open benchmark that overlapped with the original
angle. Rather than reproduce it, I narrowed the claim to the two things the released systems
still do not do. Their benchmark is also the one this dissertation evaluates on.

On a public French benchmark the fusion raises the per-span AUC from 0.814 to
0.897, a statistically significant gain obtained without fine-tuning either base model.
In domain, however, the same off-the-shelf pipeline does not transfer, and understanding
why turns out to be the more instructive result of the two.

\section{Aims and objectives}
\label{sec:objectives}

I set the project four objectives.

\begin{itemize}
\item \textbf{O1.} Build a modular, real-time endpointing pipeline (VAD, acoustic branch,
lexical branch, fusion, decision) whose decision never waits on the transcript.

\item \textbf{O2.} Fuse an off-the-shelf acoustic detector and a French language model, and
test whether the fusion beats the acoustic model alone on a public benchmark.

\item \textbf{O3.} Calibrate the fused score and choose the operating point from an
explicit cost (interruption against latency) rather than a fixed timeout.

\item \textbf{O4.} Instrument the production agent to build a weakly labelled in-domain
evaluation set, and validate the method in domain.

\end{itemize}
